\section{Coercivity of the rank 64 lift}\label{sec:lift}
Let
\[
 e_0=(2B)^{-1/2},\qquad e_n=B^{-1/2}\cos(k_nx),\quad
 k_n=n\pi/B\quad(n\ge1),
\]
be the physical even Neumann basis. Let $P_{64}$ project onto
$e_0,\ldots,e_{63}$ and $H_{64}=L+9P_{64}$. Write $H=H_{64}$ in this appendix. We prove the lift estimate
without using the stored response data.

\subsection{The entire Gamma form dominates a diagonal form}
For $a>0$ the restricted whole-line Green kernel
$e^{-a\abs{x-y}}/(2a)$ is the inverse $R_{R,a}$ of
$-\partial_x^2+a^2$ with Robin form boundary contribution
$a(\abs{v(-B)}^2+\abs{v(B)}^2)$. Removing this nonnegative term gives
the Neumann inverse $R_{N,a}$, so $R_{R,a}\preceq R_{N,a}$.
For every zero-extended local vector,
\[
 \int_0^\infty e^{-as}\norm{v(\cdot+s)-v}_2^2\dd s
 =\frac2a\norm v_2^2-2a\ip v{R_{R,a}v}
 \ge\frac2a\norm v_2^2-2a\ip v{R_{N,a}v}.
\]
Using $\rho(s)=\sum_{j\ge0}e^{-\lambda_js}$, where
$\lambda_j=2j+1/2$, and monotone summation of nonnegative forms, yields
\begin{equation}\label{eq:DN}
 t_\Gamma\succeq D_N,\qquad
 D_Ne_n=d_ne_n,\qquad
 d_n=2\sum_{j\ge0}\frac{k_n^2}{\lambda_j(\lambda_j^2+k_n^2)}.
\end{equation}
This is a form inequality on the whole local minimum domain. It does
not impose Neumann boundary conditions on an original vector. Lower
semicontinuity extends the core comparison, with inclusion in the
form domain of $D_N$.

The numbers $d_n$ are nonnegative and increasing. We give a short
rational proof that $d_{64}>9$. Set
\[
 f_t(a)=\frac{2t^2}{a(a^2+t^2)},\qquad G(t)=\sum_{j\ge0}f_t(2j+1/2).
\]
Each summand increases in $t$, and $k_{64}=80\pi>240$.
Differentiation gives
\[
 f_t''(a)=\frac{4t^2(6a^4+3a^2t^2+t^4)}{a^3(a^2+t^2)^3}>0.
\]
The mesh-two trapezoidal inequality for this convex decreasing function
therefore gives
\begin{align*}
 G(240)\ge{}&f_{240}(1/2)+f_{240}(5/2)+f_{240}(9/2)
 +\tfrac12f_{240}(13/2)\\
 &+\tfrac12\log\bigl(1+(480/13)^2\bigr),
\end{align*}
since $\int_A^\infty f_t(a)\dd a=\log(1+t^2/A^2)$.
The inequality $(1+x)^{-1}>1-x$ proves that the four rational terms
exceed, respectively,
\[
 \frac{3999}{1000},\qquad\frac{799}{1000},\qquad
 \frac{444}{1000},\qquad\frac{153}{1000}.
\]
For example they exceed $4-1/57600$, $4/5-5/57600$,
$4/9-1/6400$ and $2/13-13/115200$, respectively.
Also \eqref{eq:atanh} gives
\[
 \log2>\frac{842}{1215}>\frac{693}{1000},\qquad
 \log(15/13)>\frac17.
\]
It follows that
\[
 \tfrac12\log\bigl(1+(480/13)^2\bigr)
 >\log(480/13)=5\log2+\log(15/13)
 >\frac{693}{200}+\frac17.
\]
Consequently
\begin{equation}\label{eq:d64}
 d_{64}>G(240)>\frac{3151}{350}=9+\frac1{350}>9.
\end{equation}

\subsection{A whole operator bound for the local prime term}
Write $s=\log2$, $t=\log3$, $d=B-s$ and
$a=s/\sqrt2$, $b=s/2$, $a_3=t/\sqrt3$.
The elementary log bounds give $.69<s<.70$, $1.09<t<1.10$ and
$\log5>2B$. The lower bound for $s$ follows from $842/1215>.693$;
for $t$, the first three terms in \eqref{eq:atanh} give
$\log3>1+1/12+1/80>1.09$.
Under the unitary fold on $(0,B)$ partition that interval
(up to endpoints) into
\[
 C=(0,d),\quad L_-=(s-d,s),\quad L_+=(s,s+d),\quad
 J_0=(d,s-d).
\]
Prime 2 connects $y\in C$ to $s-y$ and $s+y$, with weight $a$,
and reflects $J_0$. Prime 4 connects the two leaves $L_-,L_+$
with weight $b$. Prime 3 reflects $S=(t-B,B)$ with weight $a_3$.
The inequalities
$2s>t$, $t+s>2B$, $t+2B<4s$ show that the leaves lie in $S$,
that their prime-3 images lie in $J_0$, and that $C$ does not meet $S$.

Define the positive number
\[
 R_p=\frac{b+a_3+\sqrt{(b+a_3)^2+8a^2}}2,
 \qquad\eta_0=2a/R_p,
 \qquad w_0=\eta_0 1_C+1_{J_0\cup L_-\cup L_+}.
\]
The weighted row ratio $(\PP_Iw_0)/w_0$ equals $R_p$ on the center
and leaves and is at most $a+a_3$ on $J_0$. Here
$\PP_I$ is the complete local prime operator in the definition of $L$.
One has $R_p>5/4$: its defining quadratic at $5/4$ is negative using
$b>1/3$, $a_3>4/7$, $2a^2>4/9$.
Moreover $a<1/2$, $a_3<121/190$ and hence $a+a_3<5/4$.
Every edge map is measure-preserving and every edge weight is
nonnegative. Applying
$2\abs{uv}\le(w_y/w_x)u^2+(w_x/w_y)v^2$ on each symmetric edge,
and summing, proves
$\abs{\ip v{\PP_Iv}}\le R_p\norm v_2^2$.
This retains all low/high interactions of the full operator.
Finally
\[
 b+a_3<\frac7{20}+\frac{121}{190}=\frac{75}{76},\quad
 2a^2<\frac{49}{100},\quad
 \left(\frac{27}{20}\right)^2-\frac{75}{76}\frac{27}{20}
 -\frac{49}{100}=\frac1{3800}>0.
\]
Thus
\begin{equation}\label{eq:primeop}
 \norm{\PP_I}<\frac{27}{20}.
\end{equation}

\subsection{The lift constant and a residual majorant}
The elementary estimate
\begin{equation}\label{eq:cgamma}
 c_\Gamma<\frac65+\frac35+\frac{11}{7}+\frac{21}{10}
 =\frac{383}{70}
\end{equation}
uses $e^{6/5}>401/125>22/7$ and $\gamma_E<3/5$.
To check the latter without evaluating Euler's constant, the harmonic
error identity gives
\[
 H_n-\log n-\gamma_E
 =\sum_{k=n}^\infty\int_k^{k+1}
       \left(\frac1x-\frac1{k+1}\right)\dd x
 >\frac1{2(n+1)}.
\]
Take $n=5$ and $\log5>8/5$ to get
$\gamma_E<H_5-8/5-1/12=3/5$.

\begin{theorem}[Rank-64 lift coercivity]\label{thm:lift}
On the complete mean-inclusive local minimum domain,
\begin{equation}\label{eq:lift}
 H_{64}=L+9P_{64}\succeq\frac{61}{28}I,
 \qquad H^{-1}\preceq\frac{28}{61}I.
\end{equation}
This is a coercivity statement for $H$, not for the unlifted $L$.
\end{theorem}
\begin{proof}
Dropping only the nonnegative pole in a lower estimate, equations
\eqref{eq:DN} and \eqref{eq:primeop} give
\[
 L\succeq D_N-(c_\Gamma+27/20)I.
\]
Below mode 64 the lifted diagonal is at least $9-c_\Gamma-27/20$;
from mode 64 onward the same bound follows from \eqref{eq:d64}.
Since $9-383/70-27/20=61/28$, \eqref{eq:lift} follows.
The pole has not been removed from the operator being inverted.
\end{proof}
For the residual contract define the larger diagonal majorant
\begin{equation}\label{eq:Dold}
 D_{\rm old}e_n=\beta_ne_n,\qquad
 \beta_n=\begin{cases}(d_n+1/2)^{-1},&n<64,\\
 (d_n-17/2)^{-1},&n\ge64.
 \end{cases}
\end{equation}
Translation Cauchy--Schwarz gives
$\norm{\PP_I}\le2(a_2+a_3+a_4)<3$, and
$c_\Gamma<11/2$, so $H\succeq D_N-17I/2+9P_{64}$.
Equation~\eqref{eq:d64} proves
$0<H^{-1}\preceq D_{\rm old}\preceq2I$.
Contract (a) uses this larger majorant. The sharper scalar bound for $H^{-1}$ is used separately in the inverse-cost estimate.

\section{An exact signed exterior prime identity}\label{sec:signed}
In this appendix write $A_B=A$.
This section is analytic and requires no stored numerical input or
exterior moment constraint. Set $g=u/\alpha$ and
$E_n=\{x:x,x+\ell_n\in O\}$. Define the nonnegative full-prime slack
\begin{equation}\label{eq:slack}
 D_p(g)=\sum_{n\ge2}a_n\int_{E_n}
 \alpha(x)\alpha(x+\ell_n)[g(x+\ell_n)-g(x)]^2\dd x.
\end{equation}
The sum includes all prime powers. The two row orientations arise from
the two diagonal terms in each edge square, as the proof below shows.

\subsection{The complete exterior row and its one reflection}
On the positive exterior fold $f=\sqrt2u$, the full prime expression is
\begin{equation}\label{eq:exteriorfold}
 (\PP_Of)(x)=\sum_{n\ge2}a_n\big[
 f(x+\ell_n)+1_{x>B+\ell_n}f(x-\ell_n)
 +1_{B<x<\ell_n-B}f(\ell_n-x)\big],\quad x>B.
\end{equation}
The first two terms are same-side orientations. The third is one
reflected opposite-side input, with one coefficient $a_n$. The cases
for $x-\ell_n$ are disjoint: positive exterior, killed interval and
negative exterior. Since $\log4<2B<\log5$, no reflection occurs for
$n=2,3,4$. The unitary fold already absorbs the factor two from the
two half-lines.

The relation $p(u,u)=\int_B^\infty f\PP_Of$ is understood as an
absolutely convergent form or weighted-duality pairing, first for
truncated sums. It is not an assertion that $\PP_Of$ belongs to
physical $L^2$ for every finite-$\wo$ vector. Equivalently it is the
$L^2$ pairing with the bounded conjugate
$M_\alpha\PP_OM_\alpha$, justified by the following row bound.
Define the nonnegative potential on all of $O$ by
\begin{equation}\label{eq:primepotential}
 V_p(x)=\alpha(x)\sum_{n\ge2}a_n
 \big[\alpha(x+\ell_n)1_O(x+\ell_n)
      +\alpha(x-\ell_n)1_O(x-\ell_n)\big].
\end{equation}
For $x>B$, either existing same-side product is at most
$A_Bn^{-25/2}$: the sum of the two distances from $B$ is at least
$\ell_n$. On the reflected branch it is at most
$A_Be^{-(25/2)(\ell_n-2B)}$, because those distances sum exactly to
$\ell_n-2B$. There is only one such branch in a row.

The two same-side budgets together satisfy
\[
 2\sum_{n\ge2}a_nn^{-25/2}\le2\sum_{n\ge2}n^{-12}
 \le2\left(2^{-12}+\int_2^\infty t^{-12}\dd t\right)
 =\frac{13}{22528}<\frac1{1000}.
\]
At $n=5$, equations \eqref{eq:coefficientcaps} and
$e^{-(25/2)(\log5-2B)}<9/10$ give cost less than $261/400$.
There is no $n=6$ coefficient. Every remaining reflection is covered by
\eqref{eq:bothtails}, since
$a_ne^{-(25/2)(\log n-2B)}\le(5/n)^{12}$ for $n\ge7$.
Consequently, for all $x\in O$,
\begin{equation}\label{eq:primecap}
 V_p(x)<\left(\frac1{1000}+\frac{261}{400}+\frac3{100}\right)A_B
 =\frac{1367}{2000}A_B<\frac7{10}A_B.
\end{equation}
Evenness extends the pointwise bound to the negative half-line without
doubling it. The summable majorants give uniform row convergence and
boundedness of the conjugated operator, including every prime tail.

\subsection{The signed identity and the analytic cap}
\begin{theorem}[Full-prime slack and mass bound]\label{thm:signed}
For every real even exterior $u$ of finite $\wo$, one has
\begin{equation}\label{eq:slackidentity}
 p(u,u)+D_p(u/\alpha)
 =\int_O V_p(x)|u(x)/\alpha(x)|^2\dd x,
\end{equation}
with $0\le D_p(u/\alpha)\le(7/5)A_B\wo(u)$.
Put
\[
 V(x)=c_\Gamma\alpha(x)^2+V_p(x),\qquad N=\frac{31}{46080}.
\]
Then $0<V(x)<N$ and
\begin{equation}\label{eq:negativeanalytic}
 c_\Gamma\norm u_2^2+p(u,u)
 \le N \wo(u)-D_p(u/\alpha).
\end{equation}
Whenever $t_\Gamma(u)$ is finite, the exact spatial identity is
\begin{equation}\label{eq:signedq0}
 q_0(u)=t_\Gamma(u)+D_p(u/\alpha)
              -\int_OV(x)|u(x)/\alpha(x)|^2\dd x.
\end{equation}
The cap is an analytic conclusion, not a numerical premise.
\end{theorem}
\begin{proof}
For a finite sum, expansion of one edge square in \eqref{eq:slack}
gives cross term $-2a_nu(x)u(x+\ell_n)$. The two diagonal terms,
after a change of variable in one, give the two orientations in
\eqref{eq:primepotential}. This proves \eqref{eq:slackidentity} for
truncated sums with exactly the factor in the definition of $p$.
The inequality $(a-b)^2\le2(a^2+b^2)$ gives
\[
 0\le D_p(g)\le2\int_OV_pg^2\le\frac75 A_B\norm g_2^2.
\]
The same row bound controls the absolute prime bilinear form by
weighted Cauchy--Schwarz. Monotone convergence of the slack series and
absolute convergence of the prime pairing extend the identity to all
$g\in L^2(O)$. Thus negative prime correlation increases the retained
slack; its sign has not been replaced by an absolute correlation.

Equation~\eqref{eq:cgamma} gives $c_\Gamma<11/2$. Together with
\eqref{eq:envelope} and \eqref{eq:primecap},
\[
 0<V(x)<\left(\frac{11}{2}+\frac7{10}\right)A_B
 =\frac{31}{5}A_B<\frac{31}{46080}.
\]
Adding the mass term proves \eqref{eq:negativeanalytic}, and
subtracting the exact identity from $t_\Gamma$ proves
\eqref{eq:signedq0}. The slack and spatial potential are $\wo$-bounded,
so the same identities extend in the original form topology, in
particular on the fixed minimum domain $\D_O^0$.
\end{proof}

The slack has trivial kernel on the even weighted exterior space.
Indeed its $n=2$ same-side edges have strictly positive weight, so
zero slack forces $g(x+\log2)=g(x)$ almost everywhere on $(B,\infty)$.
The $L^2$ norm is then an infinite sum of equal nonnegative period
integrals and can be finite only if each vanishes. Evenness supplies
the other half-line. This trivial-kernel statement does not assert a
uniform inequality $D_p\ge\delta\wo$ with $\delta>0$.

